\chapter{Дүгнэлт}
\label{ch:conclusion}

%======================================================================
% 2026.11.09: сургуулийн 5 бүлгийн бүтцээр (өмнөх ШМТК-ийн дипломын хэв
% маяг): гол үр дүн → практик ач холбогдол → хязгаарлалт → цаашдын ажил.
% 11/07-ны «Хэлэлцүүлэг» бүлгийн хязгаарлалтыг энд, тайлбарыг 4.3-т шилжүүлэв
% (Chapter5_Discussion.tex файл хадгалагдсан, main.tex-д орохгүй).
%======================================================================

\section{Судалгааны гол үр дүн}
\label{sec:concl-main}

Энэ ажилд сургуулийн ресепшнд зориулсан, Монгол хэлээр дуугаар болон бичгээр асуухад сургуулийн албан эхээс хариулдаг туслах системийн прототип хийсэн. Системийн хоёр хэсгийг тусад нь туршсан: дуу таних хэсгийг нэг хүний бичлэгээр, хайлт ба хариултын хэсгийг бичгээр асуусан асуултаар. Дуугаар асуугаад хариулт авах бүтэн урсгал, ресепшний хүснэгтийн хайлтыг энэ ноорогт хэмжээгүй.

\textbf{RQ1.} Монгол хэлээр сургасан Whisper small загвар ба Gemini 3.5 Transcribe үйлчилгээг нэг хүний 12 бичлэг дээр харьцуулсан. Whisper 38.07\%, Gemini 40.34\% үгийн алдаатай гарсан. Бичлэг цөөн тул энэ ялгаа тодорхой биш. Gemini-ийн тэмдэгтийн алдаа их гарсан нь гол төлөв оныг цифрээр бичсэнээс болсон, гэхдээ нэг бичлэгт оныг буруу таньсан. Иймээс энэ хэмжилтээр Gemini-г илүү гэж үзээгүй. Системд Gemini-г зориулалт, зардал, нэр томьёо зааж өгөх боломжоор нь сонгосон.

\textbf{RQ2.} Сургуулийн сайтын мэдээллээс BM25 ба E5-аар хайж, дүрмээр эх сонгодог системийг 40 асуултаар шалгасан. Эхтэй 20 асуултын бүгдийн зөв хэсэг хайлтын эхний 3-т орсон. Гэхдээ хариулт 13-т нь л бүрэн гарсан. Эх нь байхгүй 19 асуултын 7-д систем эхгүйгээр хариулсан. Иймээс сул тал нь хайлт биш, хариулт сонгох алхам ба эхгүй асуултыг танихад байна. Шалгасан гурван өөрчлөлт бүгд урьдчилсан дүрмээр тэнцээгүй.

Ресепшний асуултын ихэнхэд одоогийн сан хариулт агуулаагүй. Энэ нь системийн гол бэрхшээл хайлтын арга биш, өрөө, багш, хуваарийн мэдээлэл дутуу байгааг харуулсан.

\section{Судалгааны практик ач холбогдол}
\label{sec:concl-practical}

\begin{enumerate}
    \item Сургуулийн сайтын мэдээллээр бичгээр асуухад эхийн үгээр хариулдаг прототип ажиллаж байна. Хариулт бүрт эх хавсаргадаг. Гэхдээ эх нь байхгүй 19 асуултын 7-д ч систем эх хавсаргаж хариулсан тул эх харуулах нь зөв хариултын баталгаа биш. Хэрэглэгч эхийг шалгадаг эсэх, ажилтны ачаалал буурах эсэхийг хэмжээгүй.
    \item Мэдлэгийн мод сургуулийн ямар мэдээлэл дутуу байгааг тодорхой харуулсан. Үүгээр сургалтын албанаас мэдээлэл хүсэх жагсаалт гаргасан.
    \item Сургалтын албаны хүснэгтийг шалгаж оруулах програм бэлэн. Мэдээлэл ирэхэд системийг дахин сургахгүй, зөвхөн хүснэгтийг шинэчилнэ.
\end{enumerate}

\section{Хязгаарлалт}
\label{sec:limits}

\textbf{Туршилт 1 (дуу таних):}
\begin{enumerate}
    \item \textbf{Нэг хүн.} Бүх бичлэг нэг хүнийх. Өөр хүний хоолой, аялга, ресепшний чимээтэй орчинд үр дүн өөр байж болно.
    \item \textbf{Цөөн бичлэг.} 12 бичлэг, 176 үг. Интервал өргөн бөгөөд бодит тодорхойгүй байдал үүнээс ч их байж болно (\ref{sec:stt-ci}).
    \item \textbf{Хөгжүүлэлтийн багц.} Бичлэгүүдийг өмнө нь системийг хөгжүүлэхэд харсан. Тусдаа шалгалтын багц биш.
    \item \textbf{Нэг удаагийн хэмжилт.} Gemini-д бичлэг бүрийг нэг удаа илгээсэн. Нэг бичлэг хязгаарын алдаагаар хасагдсан.
    \item \textbf{Үйлчилгээний хувилбар.} Google загвараа шинэчилж болно. Энэ үр дүн 2026.10.30-ны үйлчилгээнд л хамаарна.
    \item \textbf{Тооны дүрэм.} Цифр ба үгээр бичсэн тоог ялгаатай гэж тооцсон нь Gemini-д хатуу байж болно.
\end{enumerate}

\textbf{Туршилт 2 (хайлт ба хариулт):}
\begin{enumerate}
    \item \textbf{Цөөн асуулт.} 40 асуулт, сайтын нэг хуулбар, нэг өдрийн хэмжилт. Статистикийн ялгааг батлаагүй.
    \item \textbf{Асуултыг хүн хянаагүй.} Асуултыг AI туслах бичиж, сайтын тексттэй тулгаж түгжсэн. Хүн тусад нь хянаагүй.
    \item \textbf{«Утга гарсан» нь бүрэн зөв гэсэн үг биш.} Хариултад шаардсан утга байгаа эсэхийг шалгасан. Хариултыг хүн уншиж үнэлээгүй.
    \item \textbf{Бичгийн асуулт.} Дуугаар асуухад яриа танилтын алдаа нэмэгдэнэ \cite{sidiropoulos2022asr}. Үүнийг энэ туршилтад хэмжээгүй.
    \item \textbf{Ресепшний мэдээлэлгүй.} Сургалтын албаны хүснэгт туршилтын үед ороогүй тул ресепшний хүснэгтийн хайлтыг хэмжээгүй.
\end{enumerate}

\section{Цаашдын ажил}
\label{sec:future}

\begin{enumerate}
    \item Сургалтын албаны өрөө, багш, анги, хуваарийн мэдээллийг шалгаж оруулж, мэдлэгийн модны хоосон хэсгийг бөглөх.
    \item Ресепшний асуултаар шинэ, хүн хянасан багц бэлдэж хэмжих.
    \item Өөр 2--3 хүний зөвшөөрөлтэй бичлэгээр, ресепшний бодит орчинд дуу танилтыг дахин хэмжих.
    \item Хариулт сонгох алхмыг сайжруулах. Асуух үгийн хувирлыг засах дүрэм (F1--F3) хоёр дахь багц дээр ахиц өгөөгүй тул ижил дүрмийг дахин санал болгохгүй. Өөр аргыг (жишээ нь эхгүй асуултыг тусад нь таних) хүн хянасан шинэ багцаар л шалгана.
    \item Хариултыг дуугаар унших хэсгийг хүнээр сонсуулж үнэлэх.
\end{enumerate}
